\chapter{Conclusion and Future Work}

\section{Summary}
% Add shared memory finding 
% Add CFAR contribution
% Dont repeat details about the field trial
% Mention the effects of validation on DS4. 

This thesis designed, implemented, and evaluated a software-defined radar tracking pipeline
for \gls{uav} detect-and-avoid under strict \gls{swap} constraints. The work was motivated
by the miniaturisation of millimetre-wave radar hardware, capable embedded
\gls{gpu} compute, and an evolving regulatory landscape that increasingly requires
non-cooperative target sensing for \gls{bvlos} UAV operations. Three gaps in the existing
literature were identified: the absence of end-to-end pipeline validation air-to-air radar 
tracking of small
moving airborne targets; the lack of low-level implementation benchmarking on embedded
GPU hardware representative of UAV platforms; and a lack of prior work that characterises
how algorithm and implementation conclusions transfer across hardware configurations.
Each gap was addressed, but not fully filled, through the experiments conducted in this thesis.

The pipeline was structured into three subsystems: a processing subsystem performing
formatting, windowing, an \gls{fft}, and coherent integration; a detection and tracking subsystem
performing \gls{cfar}, clustering monopulse angle sensing, blanking, and multi-target tracking; 
and
an aggregation layer combining data from multiple receiver quadruplets. The full pipeline
demonstrated real-time multi-target tracking on an NVIDIA Xavier AGX, meeting a 
requirement of processing four
video streams of 256 chirps at ten frames per second with end-to-end latency below
50\,ms. The results establish that a complete sensing-to-tracking pipeline is feasible within the
SWaP constraints of a small \gls{uav} of less than 25 kg.

Regarding the first gap, the pipeline was validated on synthetic datasets spanning a wide
range of noise intensities and target geometries, including multi-target scenarios with
targets at 500, 1000, and 1500\,m. End-to-end tracking performance was characterised
quantitatively across windowing, integration, \gls{cfar}, angle sensing, and tracking
algorithm choices. Taylor windows calibrated to $-50$\,dB sidelobe levels were found to
minimise \gls{cfar} false alarms without recall penalty; coherent integration was the most
effective mechanism for extending detection sensitivity, raising recall from near zero to
0.9 with 64 chirps; and CA-CFAR outperformed GO-CFAR and SO-CFAR in high-noise conditions
owing to its more stable noise floor estimate. The pipelines resolves a literature contention 
through the use of a CFAR technique that is both efficient and sensitive. A gap regarding
multi-target tracking was addressed through the validation of the monopulse angle sensing stage and
multi-target tracker, which were shown to be unbiased estimators of target bearing and elevation even at
$-10$\,dB SNR with three targets in the air simultaneously. 

The system was also deployed in a live
airborne field trial. Detection of a small drone target at close range was not achieved
due to a receiver-transmitter coupling, but the pipeline correctly geolocated high cross-section
static targets at up to 706\,m, confirming that the processing chain produces physically
meaningful outputs in the field. The data from the trial formed the basis for further analysis, 
which incorporated the noise from the real-world data into the simulated data to produce DS4. 
The validating in this data revealed differences in the real and assumed noise regimes, but also
confirmed robustness of the pipeline design to those differences. 

Regarding the second gap, an ablation study was performed to understand the contribution 
to latency made by each of the stages in the pipeline. On the embedded NVIDIA Jetsons, the 
capture, integration, aggregation, and cfar stages were the most expensive to implement performance-wise.
GPU-based implementations of windowing, coherent integration,
and \gls{cfar} were benchmarked against CPU alternatives across buffer sizes and training
window configurations on both the desktop and Xavier platforms. The profiling identified
\gls{cfar} as the dominant latency contributor in the detection and tracking subsystem 
and quantified
host-device data transfer and thread synchronisation as the primary remaining overhead. 
These findings provide the hardware-constrained implementation analysis that prior UAV DAA 
radar work lacked.

Regarding the third gap, the same algorithm and implementation choices were evaluated 
across three platforms of differing architecture, the desktop workstation, the Xavier NX
, and the Xavier AGX, to test whether conclusions drawn on one transfer to another. 
They did not. The fastest CFAR implementation differed by platform: on the Xavier the 
GPU implementation was preferable across all practical training-cell counts, whereas on 
the desktop the prefix-sum CPU implementation was faster until buffers exceeded roughly 
fifty chirp records, confirming that implementation selection cannot be made platform-
agnostically. More strikingly, the two unified-memory Jetson devices achieved higher 
pipeline throughput than the discrete-GPU desktop despite far lower compute. This unexpected outcome
can be attributed to their shared CPU-GPU memory removing the host-device transfer that 
dominates the desktop on this transfer-bound workload. This result inverts the 
intuition that greater GPU capability yields greater throughput, and demonstrates that 
hardware architecture changes which algorithms and implementations are viable. 
Together, the three bodies of work advance the state of 
the art from isolated sub-system demonstrations toward a unified, empirically grounded 
understanding of the design space for compact airborne radar.

\section{Significance}
% A section that links back to the significance claims in the introduction. Relate the 
% findings of the work to the significance
% \section{Significance}

The engineering significance of this work is that it lowers a specific barrier to 
putting radar on small uncrewed aircraft. Radar is the one detect-and-avoid sensor that 
keeps working through darkness, rain, dust, and glare, but it has historically been too 
heavy, too power-hungry, and too coarse in angle to carry on a sub-25,kg platform. This 
thesis shows that a complete sensing-to-tracking pipeline can run in real time within 
that platform's size, weight, and power budget, and it replaces guesswork with measured 
evidence about which processing choices make such a system sensitive enough and fast 
enough to be useful. Just as importantly, it produces a reusable way of answering that 
question, an evaluation framework of profiling, ablation, and parameter sweeps, so that 
as embedded hardware continues to change, the conclusions can be regenerated rather 
than re-guessed. This matters because the pace of the field is set less by new radar 
physics than by what can be made to run on the compute a drone can actually carry.

The larger reason this matters lies beyond the engineering. Reliable, all-weather 
sensing of non-cooperative aircraft is the single capability standing between where 
uncrewed aviation is now, confined largely to line-of-sight and segregated airspace, 
and where regulators are preparing for it to be: routine beyond-visual-line-of-sight 
flight in shared skies. The regulatory bodies now drafting the rules for that 
transition are doing without at complete understanding of the sensing capabilities 
that will be available in the future, and 
work that reduces the uncertainty about what compact radar can actually achieve feeds 
directly into that live policy process. The economic stake is substantial: the uncrewed 
area is the fastest-growing part of the airborne radar market, and the services it 
unlocks, from same-day medical and parcel logistics to continuous inspection of power 
lines, and transport corridors, depend on small aircraft being trusted to see and 
avoid hazards without a human in the loop.

In Australia, beyond-visual-line-of-sight drones are already identified as tools for 
bushfire 
monitoring and response, where the ability to fly through smoke and low visibility, 
exactly the conditions that blind optical sensors, determines whether the aircraft is 
useful at all. The same all-weather sensing that lets a drone hold separation from 
other traffic is what lets it work over a fire front or a disaster zone at night. By 
making that perception achievable within the constraints of a small, affordable platform
, radar-based detect-and-avoid helps extend these capabilities to the regional councils
, volunteer services, and small operators for whom a lightweight, low-cost sensor is 
the only viable option.

\section{Future Work}

Several directions remain open. On the algorithm side, the monopulse angle sensing stage
should be extended to handle multiple targets per range bin through target separation
strategies. The \gls{cfar} stage should be compared against a broader set of detectors,
including OS-CFAR, MTI filters, and machine learning based detectors, evaluated under
the non-Gaussian noise statistics of real airborne data. A wider sweep of clustering
algorithms including DBSCAN should be assessed with quantitative metrics and their
downstream effects on tracking performance characterised. The tracking stage should be
extended to compare GNN-Kalman against MHT, JPDAF, and PHD-based approaches. It would 
be insightful for a future study to take the per stage analysis further than what 
was accomplished in this thesis, by looking at joint algorithms such as the joint
detection and angle sensing explored in \cite{an2024_time_varying_angle_estimation_unresolved_extended_targets_monopulse}
and track before detect approaches such as particle filters. These joint approaches 
represent a departure from the pre-stage breakdowns of this thesis, and may yield
higher sensitivities at the cost of higher computational complexity. 

On the implementation and hardware side, the dominant source of remaining latency is 
host-device data transfer and thread synchronisation overhead. A valuable direction 
would be a low-level study of how to exploit unified memory in radar pipelines, 
building on this thesis's finding that eliminating host-device transfer allowed the 
embedded Jetson devices to out-throughput a far more powerful discrete-GPU workstation; 
the integration stage's current single-threaded reduction, for instance, could be 
redistributed across GPU threads under a unified-memory access pattern. The pipeline 
was evaluated on the Xavier NX and Xavier AGX; extending this evaluation to newer 
platforms such as the Orin Nano and Orin NX would provide updated hardware-aware 
recommendations as the embedded compute landscape continues to mature.

A further direction is the fusion of the radar pipeline with a co-located electro-
optical camera. The most consequential limitation of the angle-sensing stage identified 
in this work, the collapse of multiple co-range targets into a single power-weighted 
bearing and elevation estimate, stems from the large beam width of the radar. 
An optical sensor is complementary in precisely this respect: 
it provides fine angular resolution but no direct range, whereas the radar provides 
robust range and radial velocity but coarse angle. Fusing the two would allow the 
camera to resolve targets the radar cannot separate in angle while the radar continues 
to range and track through the poor-visibility conditions in which vision degrades. 
Lies et al. \cite{Lies2020} already delegate angle sensing to a camera to free radar 
compute budget for sensitivity recovery, but quantify the resulting trade-off only 
asymptotically rather than measuring it on hardware. A fusion-oriented thesis would 
extend this from a fixed division of labour to a genuinely fused estimator, and would 
treat the fusion architecture itself, whether association is performed at the track 
level, the detection level, or through a learned joint representation, as a design 
variable in the same manner that this thesis treated the choice of window, detector, 
and integration mode. 

Ultimately, the pipeline should be re-evaluated on field-obtained data with the 
transmitter-receiver coupling artefact characterised or corrected. Achieving airborne 
tracking of multiple small, non-cooperative targets in a live flight test remains the 
long-term goal for this line of work.